\chapter{逆引き：やりたいこと・起きたことから探す}

「何をしたいか」「何が起きたか」から Step を引くための索引です。

\section*{準備でやりたいこと}
\addcontentsline{toc}{section}{準備でやりたいこと}

\begin{center}
\small
\begin{longtable}{@{}P{86mm}P{22mm}P{34mm}@{}}
\toprule
\hd{やりたいこと} & \hd{Step} & \hd{補足} \\
\midrule
\endhead
コースデータを OCAD から取り込みたい & Step 5 & \\
OCAD の番号と実際のユニット番号が違う & Step 5 & 変換用ファイルを使う \\
コース数が多く，画面入力が大変 & Step 7 & \texttt{.dat} を直接編集 \\
道路横断の区間を記録から外したい & Step 6 & 番号の頭に \texttt{*} \\
クラスとコースを一括で対応づけたい & Step 6 & 一括入力ボタン \\
競技時間・入賞人数を入れたい & Step 6 & 入賞者リストの印刷に使う \\
競技者登録番号の空欄を埋めたい & Step 9 & \tRANK \\
ランキングで組分けしたい & Step 9・12 & \tRANK\ 種別で取得先が違う \\
ゼッケン番号に意味を持たせたい & Step 10 & \\
同じ所属が連続しないようにしたい & Step 12 & \\
救護・スタート・フィニッシュに配る紙を作りたい & Step 13 & \\
レンタルカードを一括で割り当てたい & Step 14 & 読み取り機で自動入力 \\
本番前に操作を練習したい & Step 16 & 再現テストツール \\
\bottomrule
\end{longtable}
\end{center}

\section*{当日やりたいこと}
\addcontentsline{toc}{section}{当日やりたいこと}

\begin{center}
\small
\begin{longtable}{@{}P{86mm}P{22mm}P{34mm}@{}}
\toprule
\hd{やりたいこと} & \hd{Step} & \hd{補足} \\
\midrule
\endhead
複数の PC で作業したい & Step 23 & ポート 2035 \\
スタート地区と情報を共有したい & Step 23 & クラウド \\
PC どうしの時計を合わせたい & Step 24 & ネットワーク時計 \\
レンタルカードだけ音を変えたい & Step 30 & 読み取り音設定 \\
ペナのときだけ紙を出したい & Step 33 & ペナレポート自動印刷 \\
速報の書式を変えたい & Step 33 & スタイルシート \\
ユニット番号を競技中に変えたい & Step 29 & コース設定変更 \\
パンチングフィニッシュの設定を忘れた & Step 29 & 同上。読み取り済みも再計算される \\
出走したか欠席かを確定したい & Step 36 & \tSI\ チェックステーションから自動判定 \\
欠席を速く入力したい & Step 36 & 【入力】→【欠席】で番号＋Enter \\
表彰対象を報告したい & Step 35 & ふりがなを添える \\
\bottomrule
\end{longtable}
\end{center}

\section*{起きたことから探す}
\addcontentsline{toc}{section}{起きたことから探す}

\begin{center}
\small
\begin{longtable}{@{}P{86mm}P{22mm}P{34mm}@{}}
\toprule
\hd{起きたこと} & \hd{Step} & \hd{補足} \\
\midrule
\endhead
CSV をドロップしても何も起きない & Step 5 & \textbf{文字コード}（UTF-8 BOM 付き） \\
クラス名・コース名が文字化けする & Step 5 & 同上 \\
カードが 1 枚だけ読めない & Step 32・37 & バックアップ計時 \\
カードが 1 枚も読めない & Step 37 & 予備の読み取り機へ \\
読んだら他人の名前が出る & Step 32・37 & スタートリストの行ずれ \\
「記録を計算できません」 & Step 29・37 & パンチングフィニッシュ未設定 \\
全員がミスパンチになる & Step 29・37 & クラスとコースの対応 \\
判定は OK だがラップが不自然 & Step 31 & フィニッシュ後のパンチ／クリア忘れ \\
別のコースとして判定されている & Step 31 & 自動判定はペナ最小のコースを選ぶ \\
未登録カードと出る & Step 32 & 当日申込の入力漏れ \\
「既存データあり」と出る & Step 32 & 選択を間違えると前の記録が消える \\
使用者を選ぶ画面が出る & Step 32 & カードの使い回し \\
ナンバーを間違えて登録した & Step 32 & 特別対応→読み取り解除 \\
\tEMIT\ NO ACTIVATE と出る & Step 32 & 通過記録を見て OK/DISQ を決める \\
\tEMIT\ S--1 のラップと総タイムがおかしい & Step 37 & スタートユニットの故障 \\
遅刻・早発の警告が出る & Step 32 & 複数人同じずれ幅ならカード取り違え \\
タイムが全体的に数分ずれる & Step 37 & \tEMIT\ スタートユニット／\tSI\ 時計ずれ \\
\tSI\ コントロールのデータが消えている & Step 37 & クリア設定になっていた可能性 \\
\tSIAC\ タッチフリーで反応しない & Step 37 & 差し込みパンチなら競技可能 \\
クラウドが「オフライン」 & Step 23 & 紙運用に切り替え \\
サブ PC がつながらない & Step 23 & IP／ポート／ファイアウォール \\
棄権した人のカードだけ届いた & Step 32 & 読まないと未帰還に残る \\
未帰還者が減らない & Step 36 & DNS と救護所棄権を先に処理 \\
Lap Center にアップロードできない & Step 39 & 紛らわしい文字の打ち直し \\
\bottomrule
\end{longtable}
\end{center}

\section*{システム別に確認したいこと}
\addcontentsline{toc}{section}{システム別に確認したいこと}

\begin{center}
\small
\begin{longtable}{@{}P{26mm}P{74mm}P{40mm}@{}}
\toprule
\hd{システム} & \hd{その日だけ必要になること} & \hd{Step} \\
\midrule
\endhead
\tEMIT & スタートユニットを渡す・動作を確認する & Step 8・22・25 \\
\tEMIT & リーディングユニットの予備（MTR でも可）を用意する & Step 2 \\
\tEMIT & バックアップラベルがカードに入っているか確認する & Step 18 \\
\tEMIT & リフトアップスタートなら遅刻を記録加算で入れる & Step 32 \\
\tEMIT & NO ACTIVATE の対応 & Step 32 \\
\midrule
\tSI & ステーションの時計を当日朝に合わせる & Step 24 \\
\tSI & 動作時間を計算して設定する & Step 8 \\
\tSI & バックアップメモリを消去する & Step 8 \\
\tSI & ピンパンチ対応表を作る & Step 18 \\
\tSI & チェックステーションから出走者を確定する & Step 36 \\
\midrule
\tSIAC & 動作時間を 12 時間にする & Step 8 \\
\tSIAC & クリアのコード番号を 1 にする & Step 8 \\
\tSIAC & 当日朝にバッテリーテストをする & Step 25 \\
\tSIAC & SIAC TEST を会場とスタートに置く & Step 2・25 \\
\tSIAC & 撤収後すぐスタンバイに戻す（電力消費が大きい） & Step 42 \\
\tSIAC & コースがフィニッシュのそばを通らないようにする & 第 0 部・第 5 部 \\
\bottomrule
\end{longtable}
\end{center}
